\documentclass[12pt, titlepage]{article}
%\usepackage[norsk]{babel}
\usepackage[utf8]{inputenc}
\usepackage[T1]{fontenc}
\usepackage{minted, enumerate}
\usepackage[colorlinks = true,
linkcolor = blue,
urlcolor = blue]{hyperref}


\begin{document}

\title{Innlevering: Mappe}
\author{IIRA2001 - Høst 2024}
\date{}
\maketitle

Mappen består av de to mappeoppgavene vi har jobbet med i høst.

\section*{Del I}

I første del av mappen skulle dere skrive 3 pythonprogrammer:
\begin{itemize}
   \item {Et som bruker den logistiske likningen, feks ved populasjonsvekst}
   \item {En sparekalkulator, \textit{eller} et program rundt databehandling av en kundegruppe}
   \item {En simulasjon av kundeadferd i butikk \textit{eller} simulasjon av markedsdynamikk}
\end{itemize}

Se egen oppgavebeskrivelse på blackboard for mer detaljer om programmene

\section*{Del II}

I andre del av mappen skulle dere lage 2 programmer som gjorde dataanalyse, brukte et web-api eller begge deler.
Vi krevde at minst 1 av programmene brukte data fra SSB eller Eurostat. Vi ser gjerne at dere bruker et web-api eller et bibliotek til å hente dataene gjennom python i 1 (begge er også greit) av programmene.

Se egen oppgavebeskrivelse på blackboard for mer detaljer om programmene


\section*{Rapport}

\begin{itemize}
   \item{Til hvert pythonprogram skulle dere også å skrive en liten rapport}
   \item {Rapporten må inneholde en beskrivelse av hva målet med programmet er, feks hvilken problemstilling du undersøker.}
   \item {Rapport og kodekommentar må tilsammen gjøre rede for hvordan koden din fungerer}
   \item {I tillegg trenger du en diskusjon eller analyse av resultatet av programmet}
\end{itemize}

Siste punkt er spesielt viktig og relevant for del 2 av mappen.

Kort fortalt trenger vi en liten rapport som beskriver hva målet med koden din var, og diskuterer hva resultatet av koden din ble, eller hva du fant ut.\\
Hvor lang teksten er avhenger av programmet ditt. 
Dersom du har laget en sparekalkulator og bruker den til å finne ut hvor mye du må spare per måned for å kjøpe deg en lamborghini i 30-årsgave til deg selv, er det kanskje færre vurderinger og diskusjoner å skrive enn om du undersøker effekten strømprisen i europa har på aksjekursen til norsk hydro.

~\\
Dere kan  levere kode og rapport på flere måter, avhengig av hva som passer deg best:
\begin{itemize}
   \item {Som 1 *.ipynb fil (Jupyter-notebook fil) per program med kodeceller for python og tekstceller for beskrivelse/resultat/diskusjon}
   \item {Som 1 *.py fil (pythonscript) med koden din og 1 PDF-fil med beskrivelse/resultat/diskusjon, per program}
   \item {Som 1 *.py- eller *.ipynb-fil per program og en samlet rapport for hele mappen, eller 2 rapporter for del 1 og del 2 av mappen}
\end{itemize}

Det er godkjent å ha flere programmer i samme *.py- eller *.ipynb-fil, men vi ser helst at dette unngås.
Leverer du rapporten i annet format enn PDF (ie. word), risikerer du at sensor ikke klarer å lese den 

Husk også på å legge ved eventuelle data dere har hentet fra SSB eller andre kilder (om dere ikke bruker webapi):

Dere har egentlig frihet til å presentere mappen deres på den måten du synes er best, men på eget ansvar:
\begin{itemize}
   \item {Vi må kunne kjøre koden deres uten for mye styr}
   \item {Programmene trenger tekst som beskriver de, og diskuterer eventuelle problemstillinger og resultater}
\end{itemize}

\section*{KI-deklarasjonsskjema}

Til mappen skal det legges ved signert og utfylt \href{https://i.ntnu.no/documents/portlet_file_entry/1305837853/NO_KI+-+deklareringsskjema+per+v%C3%A5r+2024.pdf/3b07aee3-d348-4360-e595-057ec480f1a3}{KI-deklarasjonsskjema}

Det er regnet som juks å generere besvarelse og kode ved hjelp av kunstig intelligens og levere den helt eller delvis som egen besvarelse.
Vær nøye med å redegjøre for deler av kode og besvarelse du ikke har tilvirket selv.

\section*{Inspera}

Det åpnes for innlevering på Inspera 9.12 klokken 09:00, og leveringsfristen er den 13.12 klokken 12:00 (ikke midnatt)

Hele mappen med kode, rapport og KI-deklarasjonsskjema pakkes sammen til en zip-fil og lastes opp i inspera.

Vi legger ut litt info og lenker på blackboard om hvordan du pakker sammen filer til et zip-arkiv, og hvor dere finner Jupyter-notebook filene deres


\end{document}
